\documentclass[12pt]{article}

\usepackage{fontspec}
\newfontfamily\handwriting{a-casual-handwritten-pen}[Path=../assets/fonts/,Extension=.ttf,Scale=1.5]
\newfontfamily\comic{Comic Neue}

\usepackage[utf8]{inputenc}
\usepackage{amsmath}
\usepackage{amssymb}
\usepackage{enumitem}
\usepackage{xcolor}
\usepackage{soul}
\usepackage{tikz}
\usepackage[most]{tcolorbox}
\usepackage{titlesec}
\usepackage{multicol}
\usepackage{environ}

\usetikzlibrary{fit, positioning, arrows.meta, decorations.pathreplacing}

% --- Custom Colors ---
\definecolor{primaryblue}{RGB}{42, 111, 151}
\definecolor{exbg}{RGB}{255, 240, 245}
\definecolor{rulebg}{RGB}{232, 245, 233}
\definecolor{highlight}{RGB}{255, 243, 176}
\definecolor{pinkanno}{RGB}{255, 105, 180}
\definecolor{purpleanno}{RGB}{147, 112, 219}
\definecolor{solutionbg}{RGB}{245, 245, 255}

% --- Header Styling ---
\titleformat{\section}{\color{primaryblue}\normalfont\Large\bfseries}{\thesection}{1em}{}
\titleformat{\subsection}{\color{primaryblue}\normalfont\large\bfseries}{\thesubsection}{1em}{}

% --- Friendly Boxes ---
% --- Auto-numbered Example Box ---
\newcounter{examplenum}

\newtcolorbox{friendlyexbox}[1]{
    colback=exbg,
    colframe=pinkanno!50,
    arc=5pt,
    title=\textbf{#1},
    coltitle=black,
    attach title to upper,
    after title={:\quad}
}

\newenvironment{friendlyex}{%
    \refstepcounter{examplenum}%
    \begin{friendlyexbox}{Example \theexamplenum}%
}{%
    \end{friendlyexbox}%
}

\newtcolorbox{friendlyrule}[1]{
    colback=rulebg,
    colframe=primaryblue!30,
    arc=5pt,
    fonttitle=\bfseries,
    title={#1},
    coltitle=primaryblue
}

\newcommand{\funhl}[1]{\sethlcolor{highlight}\hl{#1}}

% --- Show/Hide Solutions ---
\newif\ifshowsolutions

% Uncomment the next line to show solutions.
\showsolutionstrue

\newcommand{\solutionblank}{\vspace{25ex}}

\NewEnviron{solutionbox}{
\ifshowsolutions
\begin{tcolorbox}[
    colback=solutionbg,
    colframe=purpleanno!45,
    arc=5pt,
    title={Solution},
    coltitle=primaryblue,
    fonttitle=\bfseries
]
\BODY
\end{tcolorbox}
\else
\solutionblank
\fi
}

\begin{document}
\handwriting

\begin{center}
    \Large\textbf{\color{primaryblue}Module 3 Lesson 4\\Radical Equations} \\
    \vspace{0.2cm}
    \hrule
\end{center}

\section*{Objectives}

\begin{friendlyrule}{In this lesson, we will learn how to}
\begin{itemize}
    \item solve equations involving radicals;
    \item recognize and discard extraneous solutions;
    \item apply radical equations to real-world formulas.
\end{itemize}
\end{friendlyrule}

\section*{1. Radical Equations}

\begin{friendlyrule}{Radical Equation}
A \funhl{radical equation} is an equation in which the variable appears inside a radical, such as $\sqrt[n]{ax+b}=c$. The \funhl{radicand} is the expression under the radical, and $n$ is the \funhl{index}. Recall that $\sqrt[n]{x}=x^{1/n}$.
\end{friendlyrule}

\begin{friendlyrule}{How to Solve a Radical Equation}
\begin{enumerate}
    \item Isolate a radical on one side of the equation.
    \item Raise both sides to a power that matches the radical's index (square both sides for a square root, cube both sides for a cube root, etc.).
    \item If a radical remains, repeat steps $1$--$2$.
    \item Solve the resulting equation.
    \item Check every solution in the \textbf{original} equation.
    \item Discard any \funhl{extraneous solutions}.
\end{enumerate}
\end{friendlyrule}

\begin{friendlyrule}{Extraneous Solution}
An \funhl{extraneous solution} is a value obtained algebraically (typically after squaring both sides) that does \textbf{not} satisfy the original equation. Squaring both sides can introduce solutions that were not there originally, so every solution must be checked.
\end{friendlyrule}

\section*{2. Solving Radical Equations}

\begin{friendlyex}{}
Solve $\sqrt{5-x}-3=0$.
\end{friendlyex}

\begin{solutionbox}
\[
\sqrt{5-x}=3 \implies 5-x=9 \implies x=-4.
\]

\textbf{Check:} $\sqrt{5-(-4)}-3=\sqrt{9}-3=3-3=0$. \checkmark

\[
\boxed{x=-4}
\]
\end{solutionbox}

\begin{friendlyex}{}
Solve $\sqrt{15-2x}=x$.
\end{friendlyex}

\begin{solutionbox}
\[
15-2x=x^2 \implies x^2+2x-15=0 \implies (x+5)(x-3)=0.
\]
\[
x=-5 \text{ or } x=3.
\]

\textbf{Check } $x=-5$: $\sqrt{15-2(-5)}=\sqrt{25}=5$, but $x=-5$. Since $5\neq -5$, this is \textbf{extraneous}.

\textbf{Check } $x=3$: $\sqrt{15-2(3)}=\sqrt{9}=3=x$. \checkmark

\[
\boxed{x=3}
\]
\end{solutionbox}

\begin{friendlyex}{}
Solve $\sqrt{x^2-2}-\sqrt{x+4}=0$.
\end{friendlyex}

\begin{solutionbox}
\[
\sqrt{x^2-2}=\sqrt{x+4} \implies x^2-2=x+4 \implies x^2-x-6=0.
\]
\[
(x-3)(x+2)=0.
\]

So $x=3$ or $x=-2$.

\textbf{Check } $x=3$: $\sqrt{9-2}=\sqrt7$ and $\sqrt{3+4}=\sqrt7$. \checkmark

\textbf{Check } $x=-2$: $\sqrt{4-2}=\sqrt2$ and $\sqrt{-2+4}=\sqrt2$. \checkmark

\[
\boxed{x=3,\ x=-2}
\]
\end{solutionbox}

\begin{friendlyex}{}
Solve $\sqrt{3t+1}+\sqrt{5-t}=4$.
\end{friendlyex}

\begin{solutionbox}
Isolate one radical: $\sqrt{3t+1}=4-\sqrt{5-t}$. Square both sides:
\[
3t+1=16-8\sqrt{5-t}+(5-t)=21-t-8\sqrt{5-t}.
\]

\[
3t+1-21+t=-8\sqrt{5-t} \implies 4t-20=-8\sqrt{5-t}.
\]
\[
t-5=-2\sqrt{5-t}.
\]

Multiply by $-1$: $5-t=2\sqrt{5-t}$. Let $u=\sqrt{5-t}$, so $5-t=u^2$:
\[
u^2=2u \implies u^2-2u=0 \implies u(u-2)=0.
\]
\[
u=0 \text{ or } u=2.
\]

If $u=0$: $\sqrt{5-t}=0\implies t=5$.

If $u=2$: $\sqrt{5-t}=2\implies 5-t=4\implies t=1$.

\textbf{Check } $t=5$: $\sqrt{16}+\sqrt{0}=4+0=4$. \checkmark

\textbf{Check } $t=1$: $\sqrt{4}+\sqrt{4}=2+2=4$. \checkmark

\[
\boxed{t=5,\ t=1}
\]
\end{solutionbox}

\begin{friendlyex}{}
Solve $4+\sqrt[3]{x-6}=2$.
\end{friendlyex}

\begin{solutionbox}
\[
\sqrt[3]{x-6}=-2 \implies x-6=(-2)^3=-8 \implies x=-2.
\]

\textbf{Check:} $4+\sqrt[3]{-2-6}=4+\sqrt[3]{-8}=4+(-2)=2$. \checkmark

\[
\boxed{x=-2}
\]
\end{solutionbox}

\section*{3. An Application: Skid-Mark Speed}

\begin{friendlyrule}{Skid-Mark Formula}
The speed $v$ (in ft/sec) of a car that leaves a skid mark of length $d$ (in feet) on dry concrete satisfies
\[
v=k\sqrt{d}, \qquad k\approx 1.5.
\]
\end{friendlyrule}

\begin{friendlyex}{}
\begin{enumerate}[label=(\alph*)]
    \item A car leaves a $54$-foot skid mark. Find its speed.
    \item A car braking at $110$ ft/sec ($75$ mph) leaves a skid mark. Find the length of the skid mark.
\end{enumerate}
\end{friendlyex}

\begin{solutionbox}
\textbf{(a)} 
\[
v=1.5\sqrt{54}=1.5(3\sqrt6)=4.5\sqrt6\approx 11.02 \text{ ft/sec}.
\]

\textbf{(b)} Solve $110=1.5\sqrt{d}$ for $d$:
\[
\sqrt d=\frac{110}{1.5}=\frac{220}{3} \implies d=\left(\frac{220}{3}\right)^2=\frac{48400}{9}\approx 5377.8 \text{ ft}.
\]
\end{solutionbox}

\end{document}
